% corpus_supplement: Erdős problem #723
% source: https://www.erdosproblems.com/latex/723
% fetched_at: 2026-07-14T18:48:22Z
% payload_sha256: 13d3dc5b5d67c2dcbb32be32dead261236a5df1c6c1e967d7e2872e954cbc80a
% statement/remarks/references extracted by erdos_ingest.extract_source_page;
% rendered by erdos_ingest.render_tex — do not edit
\documentclass{article}
\usepackage{amsmath, amssymb, amsthm}
\usepackage{hyperref}
\usepackage{geometry}
\geometry{margin=1in}

\title{Erd\H{o}s Problem \#723}
\date{Last updated: 2025-08-31}
\author{Source: erdosproblems.com}

\begin{document}
\maketitle

\noindent\textbf{Status:} OPEN \quad \textbf{No prize}

\noindent\textbf{Tags:} combinatorics

\medskip
\noindent\textbf{Problem Statement:}

If there is a finite projective plane of order $n$ then must $n$ be a prime power?

A finite projective plane of order $n$ is a collection of subsets of $\{1,\ldots,n^2+n+1\}$ of size $n+1$ such that every pair of elements is contained in exactly one set.

\bigskip
\noindent\textbf{Remarks:}

These always exist if $n$ is a prime power. This conjecture has been proved for $n\leq 11$, but it is open whether there exists a projective plane of order $12$.

Bruck and Ryser \cite{BrRy49} have proved that if $n\equiv 1\pmod{4}$ or $n\equiv 2\pmod{4}$ then $n$ must be the sum of two squares. For example, this rules out $n=6$ or $n=14$. The case $n=10$ was ruled out by computer search \cite{La97}.

\bigskip
\noindent\textbf{References:}

[BrRy49] Bruck, R. H. and Ryser, H. J., The nonexistence of certain finite projective planes. Canad. J. Math. (1949), 88-93.
 
 [La97] Lam, C. W. H., The search for a finite projective plane of order {$10$}
[{MR}1103185 (92b:51013)]. (1997), 335-355.

\bigskip
\noindent\small{Source: \url{https://www.erdosproblems.com/723}}
\end{document}
